% conclusion_reasonspec.tex — Conclusion for the Reasoning Speculation paper
% Included from main_reasonspec.tex via % conclusion_reasonspec.tex — Conclusion for the Reasoning Speculation paper
% Included from main_reasonspec.tex via % conclusion_reasonspec.tex — Conclusion for the Reasoning Speculation paper
% Included from main_reasonspec.tex via % conclusion_reasonspec.tex — Conclusion for the Reasoning Speculation paper
% Included from main_reasonspec.tex via \input{sections/conclusion_reasonspec}

We introduced \textbf{\fullmethod{}} (\method{}), a training-free framework that restructures test-time reasoning compute via parallel exploration and consensus-based routing.
Rather than allocating a single token budget and relying on weak single-path signals to judge difficulty, \method{} generates $K$ short exploration paths, computes a cross-path consensus score, and adaptively routes each problem through one of three resolution strategies---accept, extend, or deliberate.
The key insight underlying our approach is that multi-path consensus is a \emph{fundamentally richer} difficulty signal than any feature extractable from a single reasoning trace: it captures inter-path variation that is information-theoretically inaccessible to single-path methods.

\paragraph{Key findings.}
Our work yields both negative and positive results of independent interest.
On the negative side, we conducted a comprehensive evaluation of all major categories of single-path adaptive budget controllers---honest features, uncertainty-based routing, and MCTS-style refinement---and found that \emph{every} controller degrades accuracy by 5--7\,pp relative to a fixed budget, demonstrating a fundamental limitation of single-path difficulty estimation in the black-box setting.
These systematic negative results motivate the shift to multi-path computation.
On the positive side, we showed that the consensus signal achieves a Spearman correlation of $\rho = -0.51$ with oracle difficulty ($p < 10^{-88}$), far exceeding any single-path feature ($|\rho| \le 0.19$).
We proved that the optimal budget-allocation policy under standard regularity conditions is a threshold policy on consensus, and that the three-route design of \method{} naturally emerges from this optimality structure (Proposition~\ref{prop:threshold}).

\paragraph{Limitations.}
Several limitations of the current work deserve acknowledgment.
First, even though the $K$ exploration paths are embarrassingly parallel and wall-clock latency remains comparable to a single full-budget pass, the \emph{total token count} increases by a factor of up to $K$ for easy questions---a cost that matters when billing is per-token rather than per-second.
Second, the consensus mechanism currently operates only on final extracted answers and does not exploit agreement or divergence in intermediate reasoning steps, potentially discarding useful signal.
Third, our evaluation focuses primarily on mathematical reasoning (GSM8K); generalization to other reasoning domains---code generation, commonsense reasoning, scientific problem-solving---remains to be validated.
Finally, the hyperparameters $K$, $\theta_{\mathrm{E}}$, and $\theta_{\mathrm{M}}$ were tuned on a held-out validation split of GSM8K, and may require task-specific adjustment for new domains.

\paragraph{Future directions.}
The \method{} framework opens several avenues for future work.
A natural extension is \emph{semantic consensus}: comparing not just final answers but intermediate reasoning steps across paths, enabling finer-grained difficulty estimation and earlier detection of divergent reasoning strategies.
Second, the hand-crafted consensus function (Eq.~\ref{eq:confidence}) could be replaced by a lightweight \emph{learned consensus model} that aggregates multi-path signals with minimal overhead, potentially improving routing accuracy without sacrificing the training-free spirit for the base reasoner.
Third, integrating \method{} with \emph{process reward models}~\citep{lightman2024lets} could enable richer path evaluation during the resolve phase, combining the strengths of consensus-based routing with step-level verification.
Finally, extending parallel exploration to \emph{latent-space reasoning}---where ``thinking tokens'' are processed in a continuous representation rather than as discrete text~\citep{hao2024training}---could reduce the per-path cost of exploration and make the framework viable for settings where token budgets are tightly constrained.


We introduced \textbf{\fullmethod{}} (\method{}), a training-free framework that restructures test-time reasoning compute via parallel exploration and consensus-based routing.
Rather than allocating a single token budget and relying on weak single-path signals to judge difficulty, \method{} generates $K$ short exploration paths, computes a cross-path consensus score, and adaptively routes each problem through one of three resolution strategies---accept, extend, or deliberate.
The key insight underlying our approach is that multi-path consensus is a \emph{fundamentally richer} difficulty signal than any feature extractable from a single reasoning trace: it captures inter-path variation that is information-theoretically inaccessible to single-path methods.

\paragraph{Key findings.}
Our work yields both negative and positive results of independent interest.
On the negative side, we conducted a comprehensive evaluation of all major categories of single-path adaptive budget controllers---honest features, uncertainty-based routing, and MCTS-style refinement---and found that \emph{every} controller degrades accuracy by 5--7\,pp relative to a fixed budget, demonstrating a fundamental limitation of single-path difficulty estimation in the black-box setting.
These systematic negative results motivate the shift to multi-path computation.
On the positive side, we showed that the consensus signal achieves a Spearman correlation of $\rho = -0.51$ with oracle difficulty ($p < 10^{-88}$), far exceeding any single-path feature ($|\rho| \le 0.19$).
We proved that the optimal budget-allocation policy under standard regularity conditions is a threshold policy on consensus, and that the three-route design of \method{} naturally emerges from this optimality structure (Proposition~\ref{prop:threshold}).

\paragraph{Limitations.}
Several limitations of the current work deserve acknowledgment.
First, even though the $K$ exploration paths are embarrassingly parallel and wall-clock latency remains comparable to a single full-budget pass, the \emph{total token count} increases by a factor of up to $K$ for easy questions---a cost that matters when billing is per-token rather than per-second.
Second, the consensus mechanism currently operates only on final extracted answers and does not exploit agreement or divergence in intermediate reasoning steps, potentially discarding useful signal.
Third, our evaluation focuses primarily on mathematical reasoning (GSM8K); generalization to other reasoning domains---code generation, commonsense reasoning, scientific problem-solving---remains to be validated.
Finally, the hyperparameters $K$, $\theta_{\mathrm{E}}$, and $\theta_{\mathrm{M}}$ were tuned on a held-out validation split of GSM8K, and may require task-specific adjustment for new domains.

\paragraph{Future directions.}
The \method{} framework opens several avenues for future work.
A natural extension is \emph{semantic consensus}: comparing not just final answers but intermediate reasoning steps across paths, enabling finer-grained difficulty estimation and earlier detection of divergent reasoning strategies.
Second, the hand-crafted consensus function (Eq.~\ref{eq:confidence}) could be replaced by a lightweight \emph{learned consensus model} that aggregates multi-path signals with minimal overhead, potentially improving routing accuracy without sacrificing the training-free spirit for the base reasoner.
Third, integrating \method{} with \emph{process reward models}~\citep{lightman2024lets} could enable richer path evaluation during the resolve phase, combining the strengths of consensus-based routing with step-level verification.
Finally, extending parallel exploration to \emph{latent-space reasoning}---where ``thinking tokens'' are processed in a continuous representation rather than as discrete text~\citep{hao2024training}---could reduce the per-path cost of exploration and make the framework viable for settings where token budgets are tightly constrained.


We introduced \textbf{\fullmethod{}} (\method{}), a training-free framework that restructures test-time reasoning compute via parallel exploration and consensus-based routing.
Rather than allocating a single token budget and relying on weak single-path signals to judge difficulty, \method{} generates $K$ short exploration paths, computes a cross-path consensus score, and adaptively routes each problem through one of three resolution strategies---accept, extend, or deliberate.
The key insight underlying our approach is that multi-path consensus is a \emph{fundamentally richer} difficulty signal than any feature extractable from a single reasoning trace: it captures inter-path variation that is information-theoretically inaccessible to single-path methods.

\paragraph{Key findings.}
Our work yields both negative and positive results of independent interest.
On the negative side, we conducted a comprehensive evaluation of all major categories of single-path adaptive budget controllers---honest features, uncertainty-based routing, and MCTS-style refinement---and found that \emph{every} controller degrades accuracy by 5--7\,pp relative to a fixed budget, demonstrating a fundamental limitation of single-path difficulty estimation in the black-box setting.
These systematic negative results motivate the shift to multi-path computation.
On the positive side, we showed that the consensus signal achieves a Spearman correlation of $\rho = -0.51$ with oracle difficulty ($p < 10^{-88}$), far exceeding any single-path feature ($|\rho| \le 0.19$).
We proved that the optimal budget-allocation policy under standard regularity conditions is a threshold policy on consensus, and that the three-route design of \method{} naturally emerges from this optimality structure (Proposition~\ref{prop:threshold}).

\paragraph{Limitations.}
Several limitations of the current work deserve acknowledgment.
First, even though the $K$ exploration paths are embarrassingly parallel and wall-clock latency remains comparable to a single full-budget pass, the \emph{total token count} increases by a factor of up to $K$ for easy questions---a cost that matters when billing is per-token rather than per-second.
Second, the consensus mechanism currently operates only on final extracted answers and does not exploit agreement or divergence in intermediate reasoning steps, potentially discarding useful signal.
Third, our evaluation focuses primarily on mathematical reasoning (GSM8K); generalization to other reasoning domains---code generation, commonsense reasoning, scientific problem-solving---remains to be validated.
Finally, the hyperparameters $K$, $\theta_{\mathrm{E}}$, and $\theta_{\mathrm{M}}$ were tuned on a held-out validation split of GSM8K, and may require task-specific adjustment for new domains.

\paragraph{Future directions.}
The \method{} framework opens several avenues for future work.
A natural extension is \emph{semantic consensus}: comparing not just final answers but intermediate reasoning steps across paths, enabling finer-grained difficulty estimation and earlier detection of divergent reasoning strategies.
Second, the hand-crafted consensus function (Eq.~\ref{eq:confidence}) could be replaced by a lightweight \emph{learned consensus model} that aggregates multi-path signals with minimal overhead, potentially improving routing accuracy without sacrificing the training-free spirit for the base reasoner.
Third, integrating \method{} with \emph{process reward models}~\citep{lightman2024lets} could enable richer path evaluation during the resolve phase, combining the strengths of consensus-based routing with step-level verification.
Finally, extending parallel exploration to \emph{latent-space reasoning}---where ``thinking tokens'' are processed in a continuous representation rather than as discrete text~\citep{hao2024training}---could reduce the per-path cost of exploration and make the framework viable for settings where token budgets are tightly constrained.


We introduced \textbf{\fullmethod{}} (\method{}), a training-free framework that restructures test-time reasoning compute via parallel exploration and consensus-based routing.
Rather than allocating a single token budget and relying on weak single-path signals to judge difficulty, \method{} generates $K$ short exploration paths, computes a cross-path consensus score, and adaptively routes each problem through one of three resolution strategies---accept, extend, or deliberate.
The key insight underlying our approach is that multi-path consensus is a \emph{fundamentally richer} difficulty signal than any feature extractable from a single reasoning trace: it captures inter-path variation that is information-theoretically inaccessible to single-path methods.

\paragraph{Key findings.}
Our work yields both negative and positive results of independent interest.
On the negative side, we conducted a comprehensive evaluation of all major categories of single-path adaptive budget controllers---honest features, uncertainty-based routing, and MCTS-style refinement---and found that \emph{every} controller degrades accuracy by 5--7\,pp relative to a fixed budget, demonstrating a fundamental limitation of single-path difficulty estimation in the black-box setting.
These systematic negative results motivate the shift to multi-path computation.
On the positive side, we showed that the consensus signal achieves a Spearman correlation of $\rho = -0.51$ with oracle difficulty ($p < 10^{-88}$), far exceeding any single-path feature ($|\rho| \le 0.19$).
We proved that the optimal budget-allocation policy under standard regularity conditions is a threshold policy on consensus, and that the three-route design of \method{} naturally emerges from this optimality structure (Proposition~\ref{prop:threshold}).

\paragraph{Limitations.}
Several limitations of the current work deserve acknowledgment.
First, even though the $K$ exploration paths are embarrassingly parallel and wall-clock latency remains comparable to a single full-budget pass, the \emph{total token count} increases by a factor of up to $K$ for easy questions---a cost that matters when billing is per-token rather than per-second.
Second, the consensus mechanism currently operates only on final extracted answers and does not exploit agreement or divergence in intermediate reasoning steps, potentially discarding useful signal.
Third, our evaluation focuses primarily on mathematical reasoning (GSM8K); generalization to other reasoning domains---code generation, commonsense reasoning, scientific problem-solving---remains to be validated.
Finally, the hyperparameters $K$, $\theta_{\mathrm{E}}$, and $\theta_{\mathrm{M}}$ were tuned on a held-out validation split of GSM8K, and may require task-specific adjustment for new domains.

\paragraph{Future directions.}
The \method{} framework opens several avenues for future work.
A natural extension is \emph{semantic consensus}: comparing not just final answers but intermediate reasoning steps across paths, enabling finer-grained difficulty estimation and earlier detection of divergent reasoning strategies.
Second, the hand-crafted consensus function (Eq.~\ref{eq:confidence}) could be replaced by a lightweight \emph{learned consensus model} that aggregates multi-path signals with minimal overhead, potentially improving routing accuracy without sacrificing the training-free spirit for the base reasoner.
Third, integrating \method{} with \emph{process reward models}~\citep{lightman2024lets} could enable richer path evaluation during the resolve phase, combining the strengths of consensus-based routing with step-level verification.
Finally, extending parallel exploration to \emph{latent-space reasoning}---where ``thinking tokens'' are processed in a continuous representation rather than as discrete text~\citep{hao2024training}---could reduce the per-path cost of exploration and make the framework viable for settings where token budgets are tightly constrained.
